\documentclass[10pt,a4paper,ssfamily]{exam}
\usepackage[brazil]{babel}
\usepackage[numbers,sort&compress]{natbib}
\usepackage[utf8]{inputenc}
\renewcommand{\baselinestretch}{1.0}
\usepackage{epsfig}
\usepackage{mhchem}
\usepackage{graphicx}
\usepackage{makeidx}
\usepackage{multicol}
\usepackage{multirow}
\usepackage{amssymb}
\usepackage{amsmath}
\DeclareMathOperator{\sen}{sen}
\newcommand{\listfont}{\bf\fontencoding{OT1}\sffamily\large}
\newcommand{\titlesfont}{\fontencoding{OT1}\sffamily\large}
\renewcommand{\baselinestretch}{1.0}
\newcommand{\Mg}{Mg$^{2+}$}
\renewcommand{\t}{\textrm}
\newcommand{\cc}[1]{[\textrm{#1}]}
\usepackage{enumitem}
\setlist{nosep, topsep=2pt, partopsep=0pt, parsep=0pt, itemsep=2pt}
\newcommand{\bmat}{\left[\begin{array}{cc}}
\newcommand{\emat}{\end{array}\right]}
\newcommand{\nomera}{
  \noindent
  \begin{tabular}{|p{0.7\textwidth}|p{0.3\textwidth}|}
  Nome: & RA: \\
  \hline
  \end{tabular}
}

\newcommand{\qini}{\begin{enumerate}[resume]\item}
\newcommand{\qend}{\end{enumerate}}

\newcommand{\oC}{$^\circ$C}
\newcommand{\kcalmol}{~kcal~mol$^{-1}$}
\newcommand{\moll}{~mol~L$^{-1}$}
\newcommand{\molkg}{~mol~kg$^{-1}$}
\newcommand{\gl}{~g~L$^{-1}$}
\newcommand{\e}[1]{$\times 10^{#1}$}

\newcommand{\eps}{\varepsilon}

\newcommand{\Resposta}[1]{\vspace{0.5cm}\noindent{\bf Resposta:}\\ #1
\begin{center}\rule{\textwidth}{0.4pt}\end{center}}
\renewcommand{\Resposta}[1]{}
\newcommand{\resposta}[1]{\vspace{0.5cm}\noindent{\bf Resposta:}\\ #1
\begin{center}\rule{\textwidth}{0.4pt}\end{center}}


\begin{document}
\sffamily
\pagestyle{empty}
\nomera

\begin{center}
{\bf QG101 - Química Geral}\\
{\bf Aula 5 - Propriedades Periódicas}\\
\underline{Justifique suas respostas.}
\end{center}

\vspace{0.2cm}
\begin{center}{\bf Questões}\end{center}

% ---------------------------------------------------------------
\qini
O \textbf{sódio} (\ce{Na}, $Z = 11$) e o \textbf{cloro} (\ce{Cl},
$Z = 17$) são vizinhos no 3.º período e reagem para formar o cloreto
de sódio (\ce{NaCl}).

\begin{enumerate}
  \item Escreva as configurações eletrônicas de \ce{Na} e \ce{Cl}.
        Com base nas tendências periódicas, compare os raios atômicos,
        as energias de ionização e as eletronegatividades desses dois
        elementos.
  \item O \ce{Na} perde 1 elétron e o \ce{Cl} ganha 1 elétron ao
        formar o \ce{NaCl}. Compare os raios do \textbf{íon}
        \ce{Na+} com o átomo \ce{Na}, e do \textbf{íon} \ce{Cl-}
        com o átomo \ce{Cl}. Justifique cada comparação.
  \item A diferença de eletronegatividade entre \ce{Na} e \ce{Cl}
        é $\Delta\chi = 2{,}1$. O que isso indica sobre o tipo de
        ligação formada? Calcule $\Delta\chi$ para \ce{HCl}
        ($\chi_\text{H} = 2{,}1$; $\chi_\text{Cl} = 3{,}0$) e
        indique o tipo de ligação.
\end{enumerate}

\vspace{0.3cm}
\qend

% ---------------------------------------------------------------
\qini
A tabela abaixo apresenta energias de ionização sucessivas (em
kJ/mol) de três elementos do 3.º período: \ce{Na}, \ce{Mg} e
\ce{Al} — identificados como \textbf{X}, \textbf{Y} e \textbf{Z}
(não necessariamente nessa ordem).

\begin{center}
\begin{tabular}{lrrrrr}
\hline
Elemento & EI$_1$ & EI$_2$ & EI$_3$ & EI$_4$ & EI$_5$ \\
\hline
X & 496 & 4562 & 6912 & 9543 & 13353 \\
Y & 578 & 1817 & 2745 & 11578 & 14831 \\
Z & 738 & 1451 & 7733 & 10540 & 13630 \\
\hline
\end{tabular}
\end{center}

\begin{enumerate}
  \item Para cada elemento, identifique o número de elétrons de
        valência a partir do \textbf{salto brusco} nas energias de
        ionização. Associe cada rótulo (X, Y, Z) a um dos elementos
        (\ce{Na}, \ce{Mg}, \ce{Al}), justificando sua resposta.
  \item A EI$_1$ do \ce{Mg} é \textbf{maior} que a do \ce{Al},
        mesmo que \ce{Mg} venha antes de \ce{Al} no 3.º período.
        Explique essa anomalia com base nas configurações eletrônicas
        dos dois elementos.
  \item Escreva a configuração eletrônica de \ce{Na}, \ce{Mg} e
        \ce{Al} e confirme, para cada um, o grupo da tabela periódica
        ao qual pertence.
\end{enumerate}

\vspace{0.3cm}
\qend

% ---------------------------------------------------------------
\qini
Os \textbf{halogênios} (grupo 17) formam uma família com
propriedades bem definidas. Considere \ce{F} ($Z=9$),
\ce{Cl} ($Z=17$) e \ce{Br} ($Z=35$).

\begin{enumerate}
  \item Escreva a configuração eletrônica de valência dos três
        halogênios. Por que todos eles tendem a formar o íon
        \ce{X-}?
  \item Coloque \ce{F}, \ce{Cl} e \ce{Br} em ordem crescente de:
        (i) raio atômico; (ii) energia de ionização;
        (iii) eletronegatividade. Justifique cada ordenação.
\end{enumerate}

\vspace{0.3cm}
\qend

\end{document}
